\section{Information Access}
On a daily basis, we use our personal computer for accessing information. Search engines (Google, Bing, Yahoo!, DuckDuckGo, etc.) and web browsers (Mozilla Firefox, Google Chrome, Microsoft Edge, etc.) help us to fullfill our \emph{information needs} by proving the relevant \emph{information objects}~\cite{entityorientedsearch2018}. The former aim to find web pages that contain our information need, often represented as a \emph{query}, the latter allow us to browse web pages to finally find the required information. Web browsers and search engines are two primary tools for facilitating \emph{information access}.
In this thesis, we adopt the following definition of information access---based on \textcite{encyclopedia_information_access}---which is limited to the digital activities a user undertakes to satisfy an information need. Broader considerations related to social, economic, or political factors that may restrict individuals’ freedom to access information are beyond the scope of this work.
\begin{definition}[Information Access]
    Information access is the ability to effectively identify, retrieve, and utilize information.
\end{definition}
\todo{list the possible ways for accessing information and how documents need to be processed for the different ways}
\todo{continue with something about the types of information objects, trying to reach entities.}
\todo{read https://www.encyclopedia.com/computing/news-wires-white-papers-and-books/information-access, add definition, and cite}
% information access: you know what you want
\todo{maybe discovery is better. no access is better because users know what they want. I did not work on discovery}
% https://asistdl.onlinelibrary.wiley.com/doi/abs/10.1002/(SICI)1097-4571(1999)50:9%3C737::AID-ASI2%3E3.0.CO;2-C
% information discovery: exploring or discovering new information
\todo{use information access narrowing it down to avoid confusion with socio-economic factors}

\begin{figure}[ht!]
    \centering
    \includegraphics[width=\linewidth]{images/Screenshot Rome - Google Search.png}
    \caption{Example of a modern search engine result page (screenshot from Google Search).}
    \label{fig:rome-search}
\end{figure}

\begin{figure}[ht!]
    \centering
    \includegraphics[width=\linewidth]{images/Screenshot Barack Obama - Google Search.png}
    \caption{Example of factual question answered via knowledge graph (screenshot from Google Search).}
    \label{fig:barack-search}
\end{figure}

While traditionally the information objects were documents, search engines have, over the past decades, gradually transitioned toward providing richer answers. As visible in \Cref{fig:rome-search,fig:barack-search}, results now often display entities and facts directly, enriched with pictures, maps, or other structured information depending on the answer type. A key enabler of these developments has been the emergence of large-scale knowledge bases and knowledge repositories, which organize information around entities~\cite{entityorientedsearch2018}, such as the Google Knowledge Graph~\cite{singhal_knowledge_graph,google-kg2020}.
%
\begin{figure}[ht!]
    \centering
    \includegraphics[width=\linewidth]{images/Screenshot Better Actor - Google Search.png}
    \caption{Example of question answered with a \acl{llm} (screenshot from Google Search).}
    \label{fig:actor-search}
\end{figure}
%
More recently, \acfp{llm} have further extended these capabilities by enabling natural language questions to be answered directly with fluent natural language responses, as depicted in \Cref{fig:actor-search}.

Applications for information access include, but are not limited to:
\begin{itemize}
    \item search engines that retrieve documents or items based on a query;
    \item web browsers that allow users to explore documents and follow hyperlinks;
    \item recommendation systems that suggest relevant or similar documents;
    \item question-answering systems that provide direct answers to users' questions based on information from structured or unstructured data.
\end{itemize} 